


A BLT-set of $Q(4,q)$ is a set of $q+1$ point on the quadric 
such that no point on the quadric is collinear to more than two points of the set. 
BLT sets are related to spreads of $\PG(3,q),$ to flocks 
of the quadratic cone in $\PG(3,q)$, 
and other objects in finite geometry.
These sets have been introduced in~\cite{BaderLunardonThas90} 
for the purposes of creating new flocks from old.
A BLT-set exists if and only if $q$ is odd. 
It is an interesting problem to classify BLT-sets of $Q(4,q)$ under the orthogonal group.
Some references are 
Law~\cite{Law03}, Penttila-Royle~\cite{PenttilaRoyle98}, 
Penttila-Law~\cite{LawPenttila03,LawPenttila04}, 
Betten~\cite{Betten13}, 
AlAzemi-Betten-Chowdhury~\cite{AlAzemiBettenChowdhury18}.
For small values of $q$, computer search has been used extensively. 
As we will discuss below, Orbiter is useful to supporting such computations.

\bigskip

Before we go into the problem of classification, let us first discuss how known BLT-sets 
can be constructed in Orbiter. 
Orbiter maintains a catalogue of BLT-sets over small fields. 
Besides that, 
Orbiter can be used to create instances of known families of BLT-sets.


\bigskip

Table~\ref{tab:BLT:create} shows options to create known BLT-sets.
\orbitertableofcommands{BLT_set_create}

\bigskip

Table~\ref{tab:BLT:families} shows names for known families of sporadic sets.
\orbitertableofextracommands{BLT_set_families}

\bigskip

Table~\ref{tab:BLT:activity} shows Orbiter activities for a BLT-set object.
\orbitertableofcommands{blt_set_activity}



Table~\ref{tab:BLT:classify} shows options to prepare for the classification of BLT-sets.
\orbitertableofcommands{blt_set_classify}


Table~\ref{tab:BLT:classify:activity} shows commands for the classification of BLT-sets.
\orbitertableofcommands{blt_set_classify_activity}



\bigskip


Let us demonstrate some examples. 
The command

\sourcecode{BLT_5_1}

creates the first BLT-set of order 5 from the catalogue. 
Suppose we want to export the BLT-set to GAP/fining. 
The following command can do that:

\sourcecode{BLT_5_1_export_gap}

The command creates the following GAP/fining file:
\orbiteroutput{GRAPHICS/LATEX/BLT_5_1_fining_export}

\bigskip

The command

\sourcecode{BLT_11_0}

creates the BLT-set \#0 in $Q(4,11).$
The command

\sourcecode{BLT_11_Mondello}

creates the Mondello BLT-set in $Q(4,11).$
Orbiter creates the following report:
\orbiteroutput{GRAPHICS/LATEX/BLT_mondello}



\bigskip

The classification of BLT-sets is a difficult problem. 
One approach is by means of the poset of partial BLT-sets. 
The orbits at level $q+1$ in the poset of partial BLT-sets 
are the isomorphism classes of BLT-sets. 
We illustrate this method by a small example, comparing different construction algorithms.

We begin with a classification of the complete poset of partial BLT sets. 
This method is very limited, as the number of orbits of partial BLT-sets 
gets out of hand quickly as $q$ increases.


\sourcecode{BLT_13_deep_search}

The technique of classification via substructure can help.
Various authors have described different versions of this technique,  
see~\cite{AlAzemiBettenChowdhury18,Betten13,Law03}.
As an example, we classify the BLT-sets of order $13$ 
with the technique of substructures and rainbow cliques.

\bigskip

We consider the classification problem of BLT-sets of order 13. 
We use the technique of substructures of size 5.
So, to begin with, we classify all partial BLT-sets 
of size $5$. 
For each isomorphism type of partial BLT-set of size 5, 
we create a colored graph.
The vertices in the graph are the points in $Q(4,13)$ 
which can be added to the partial BLT-set.
Two vertices are connected if they can both be added.
Here is the command:

\sourcecode{BLT_13_classify_starter}

In the next step, we compute all rainbow cliques in each of the graphs:

\sourcecode{BLT_13_cliques}

Next, we create a data structure for isomorphism testing. 
The first step is to create a database of 
all partial BLT-sets of order at most 5:

\sourcecode{BLT_13_isomorph_read_DB}

The next step is to read the rainbow cliques from the clique finding process:

\sourcecode{BLT_13_isomorph_read_solutions}

Once the cliques have been read from file, 
we compute the flag orbits in the relation between partial BLT-sets and full BLT-sets.
For each orbit representative of a partial BLT-set, 
we compute the orbits of the stabilizer on the 
BLT-sets incident with it. The following command can be used.

\sourcecode{BLT_13_isomorph_stabilizer_orbits}

Finally, we perform isomorphism testing:

\sourcecode{BLT_13_isomorph_testing}

This last command results in exactly three 
isomorphism types of BLT-sets of order 13.
Of course, these results match with the results 
from the classification of the complete poset of partial BLT-sets before.
The difference is that the second approach is more efficient, 
as the number of 
orbits to be considered by the algorithm is much smaller.
In essence, we are skipping over the worst part of the poset. 
Many partial BLT-sets do not complete to BLT-sets, 
so these braches of the search tree will die out in the clique finding algorithm.

\bigskip


The command 

\sourcecode{blt_set_export_gap_all}

exports all BLT-sets in the Orbiter catalogue to GAP. 


\bigskip

